\setcounter{ex}{0}
\setcounter{paragraph}{0}
\PhanI
\begin{ex}%book %[3P1Y3-1][Loai1] 
(QG - 2015) Một con lắc lò xo gồm một vật nhỏ khối lượng m và lò xo có độ cứng k. Con lắc dao động điều hòa với tần số góc là
\choice
{$2\pi \sqrt{\dfrac{m}{k}}$}
{$2\pi \sqrt{\dfrac{k}{m}}$}
{$\sqrt{\dfrac{m}{k}}$}
{\True $\sqrt{\dfrac{k}{m}}$}
\loigiai{
Tần số góc của con lắc lò xo: $\omega  = \sqrt{\dfrac{k}{m}}$.
}
\end{ex} \haidong{20}
\begin{ex}%book%[3P1B3-1][Loai1]
Con lắc lò xo dao động điều hòa. Khi tăng khối lượng của vật lên 9 lần thì tần số dao động của vật
\choice
{tăng lên 9 lần}
{giảm đi 9 lần}
{tăng lên 3 lần}
{\True giảm đi 3 lần}
\loigiai{
Ta có: $f=\dfrac{1}{2\pi}\sqrt{\dfrac{k}{m}} \Rightarrow$ m tăng 9 lần thì f giảm đi 3 lần.
}
\end{ex} \haidong{20}
\begin{ex}%book %[3P1B3-1][Loai1]  
Cho con lắc lò xo đang dao động điều hòa với biên độ 3 cm, biết chu kì dao động của con lắc là 0,3 s. Nếu kích thích cho con lắc dao động điều hòa với biên độ 6 cm thì chu kì dao động của con lắc bằng
\choice
{\True 0,3 s}
{0,15 s}
{0,6 s}
{0,423 s}
\loigiai{
Do chu kì dao động của vật không phụ thuộc vào các yếu tố bên ngoài mà chỉ phụ thuộc vào đặc tính riêng của hệ nên khi thay đổi biên độ sẽ không làm thay đổi chu kì.
}
\end{ex} \haidong{20}
\begin{ex}%book%[3P1B3-1][Loai1] 
Kích thích để con lắc lò xo dao động điều hòa theo phương ngang với biên độ 5 cm thì vật dao động với tần số 5 Hz. Treo hệ lò xo trên theo phương thẳng đứng rồi kích thích để con lắc lò xo dao động điều hòa với biên độ 3 cm thì tần số dao động của vật là 
\choice
{\True 5 Hz}
{10 Hz}
{15 Hz}
{6 Hz}
\loigiai{
\begin{itemize}
\item Tần số dao động của con lắc lò xo: $f=\dfrac{1}{2\pi}\sqrt{\dfrac{k}{m}}$ không phụ thuộc vào cách treo và biên độ. 
\item Do vậy với m và k không đổi, tần số của con lắc không đổi.
\end{itemize}
}
\end{ex} \haidong{20}
\begin{ex}%book %[3P1B3-1][Loai1]  
Khi gắn quả nặng có khối lượng $m_1$ vào một lò xo, thấy nó dao động điều hòa với chu kì $T_1$. Khi gắn quả nặng có khối lượng $m_2$ vào lò xo đó, nó dao động với chu kì $T_2$. Nếu gắn đồng thời $m_1$ và $m_2$ cũng vào lò xo đó, thì chu kì dao động của chúng là
\choice
{$T=T_1^2+T_2^2$}
{$T=\dfrac{{T_1+T_2}}{2}$}
{$T=T_1+T_2$}
{\True $T=\sqrt{{T_1^2+T_2^2}}$}
\loigiai{
Ta có: $T\sim \sqrt{m}\xrightarrow{{m=m_1+m_2}}{T}=\sqrt{{T_1^2+T_2^2}}$.
}


\end{ex} \haidong{20}
%\loai{Tính tần số góc}
\begin{ex} %[3P1Y3-2][Loai1] 
(QG - 2018) Một con lắc lò xo có $k = 40\ N/m$ và $m = 100\ g$. Dao động riêng của con lắc này có tần số góc là
\choice
{$400\ rad/s$}
{$0,1\pi\ rad/s$}
{\True $20\ rad/s$}
{$0,2\pi\ rad/s$}
\loigiai{
Tần số góc: $\omega =\sqrt{\dfrac{k}{m}}=20\ rad/s$.
}
%\haidong{2}
\end{ex} \haidong{20}
%\loai{Từ tần số góc tính k hoặc m}

%\loai{Tính chu kì, tần số}
\begin{ex}%[3P1B3-2]
Một con lắc lò xo, vật nặng có khối lượng $m = 250$ g, lò xo có độ cứng $k = 100$ N/m. Tần số dao động của con lắc là
\choice
{20 Hz}
{\True 3,18 Hz}
{6,28 Hz}
{5 Hz}
\loigiai{
Tần số dao động của con lắc lò xo là: $f = \dfrac{1}{{2\pi}}\sqrt{\dfrac{k}{m}} = \dfrac{1}{{2\pi}}\sqrt {\dfrac{100}{0,25}} = 3,18\ Hz$.
}
%\haidong{2}
\end{ex} \haidong{20}
%\loai{Từ chu kì, tần số tính k hoặc m}
\begin{ex} %[3P1B3-2][Loai1] 
(QG - 2018) Một con lắc lò xo gồm vật nhỏ và lò xo nhẹ có độ cứng 10\ N/m, dao động điều hòa với chu kì riêng 1 s. Khối lượng của vật là
\choice
{100 g}
{\True 250 g}
{200 g}
{150 g}
\loigiai{
Ta có: $T=2\pi \sqrt{\dfrac{m}{k}}\Rightarrow m=\dfrac{T^2.k}{4{\pi}^2}=\dfrac{1.10}{4.{\pi}^2}=0,25\ kg=250\ g$.
}

%\haidong{2}
\end{ex} \haidong{20}

%\loai{Liên hệ hai chu kì, hai tần số, hai khối lượng}
\begin{ex}
Một con lắc lò xo gồm một lò xo có độ cứng không đổi. Khi khối lượng quả nặng là m thì tần số dao động là 1 Hz. Khi khối lượng quả nặng là 2m thì tần số là
\choice
{$2\ Hz$}
{$\sqrt 2\ Hz$}
{\True $\dfrac{1}{\sqrt 2}\ Hz$}
{$0,5\ Hz$}
\loigiai{
Công thức: $f = \dfrac{1}{2\pi}\sqrt{\dfrac{k}{m}}$
}

%\haidong{2}
\end{ex} \haidong{20}
\begin{ex}%[3P1-3.2B]
Khi gắn vật nặng có khối lượng $m_1 = 4$ kg vào một lò xo có khối lượng không đáng kể, hệ dao động điều hòa với chu kì $T_1 = 1$ s. Khi gắn một vật khác có khối lượng $m_{2}$ vào lò xo trên thì hệ dao động với chu kì $T_2 = 0,5$ s. Khối lượng $m_2$ bằng
\choice
{0,5 kg}
{2 kg}
{\True 1 kg}
{3 kg}
\loigiai{
Ta có: $\begin{cases}
1\ s = 2\pi \sqrt {\dfrac{4}{k}}. \\
0,5\ s =2\pi \sqrt {\dfrac{m}{k}}. 
\end{cases} \Rightarrow 2 = \sqrt {\dfrac{4}{m}} \Rightarrow m = 1\ kg$.
}
%\haidong{2}
\end{ex} \haidong{20}

\begin{ex}%[3P1G3-2][Loai1] 
Một con lắc lò xo dao động điều hoà với chu kì T, để chu kì dao động tăng lên 10\% thì khối lượng của vật phải
\choice
{\True tăng 21\%}
{giảm 11\%}
{giảm 10\%}
{tăng 20\%}
\loigiai{
Ta có: $\dfrac{m'}{m}=\left(\dfrac{T'}{T} \right)^2=\left(\dfrac{110}{100}
\right)^2=1,21 \Rightarrow$ khối lượng phải tăng thêm 21\%.
}%<MyLT1>
%\haidong{4}
\end{ex} \haidong{20}
\begin{ex}%[3P1B3-2]
Con lắc lò xo có khối lượng m đang dao động điều hòa với chu kì 2 s. Khi tăng khối lượng của con lắc thêm 210 g thì chu kì dao động điều hòa của nó là 2,2 s. Khối lượng m bằng
\choice
{2 kg}
{\True 1 kg}
{2,5 kg}
{1,5 kg}
\loigiai{
Ta có: $\begin{cases}2 = 2\pi \sqrt{\dfrac{m}{k}} \\2,2=2\pi \sqrt {\dfrac{m + 0,21}{k}} \end{cases} \Rightarrow \dfrac{2}{2,2} = \sqrt {\dfrac{m}{m + 0,21}} \Rightarrow m = 1\ kg.$
}

%\haidong{4}
\end{ex} \haidong{20}

%\loai{Số dao động trong thời gian $\Delta t$}
\begin{ex}%[3P1B3-2]
Một con lắc lò xo dao động điều hoà. Trong khoảng thời gian  $\Delta t$, con lắc thực hiện 60 dao động toàn phần. Thay đổi khối lượng con lắc một lượng 440 g thì cũng trong khoảng thời gian $\Delta t$ ấy, nó thực hiện 50 dao động toàn phần. Khối lượng ban đầu của con lắc là
\choice
{1,44 kg}
{0,6 kg}
{0,8 kg}
{\True 1 kg}
\loigiai{
Ta có: $\begin{cases}
\dfrac{\Delta t}{60} = 2\pi \sqrt{\dfrac{m}{k}} \\
\dfrac{\Delta t}{50}=2\pi \sqrt{\dfrac{m\pm 0,44}{k}} 
\end{cases} 
\Rightarrow \dfrac{5}{6} = \sqrt{\dfrac{m}{m \pm 0,44}} \Rightarrow \dfrac{5}{6} = \sqrt{\dfrac{m}{m + 0,44}} \Rightarrow m = 1\ kg.$
}
%\haidong{6}
\end{ex} \haidong{20}

%\loai{Chu kì, tần số của vật có khối lượng $m=a.m_1+b.m_2$}


\begin{ex}%[3P1K3-2]
Khi gắn quả cầu khối lượng $m_{1}$ vào lò xo thì nó dao động với chu kì $T_1$. Khi gắn quả cầu có khối lượng $m_{2}$ vào lò xo trên thì nó dao động với chu kì $T_2 = 0,4$ s. Nếu gắn đồng thời hai quả cầu vào lò xo thì nó dao động với chu kì $T = 0,5$ s. Giá trị $T_1$ là
\choice
{0,45 s}
{\True 0,3 s}
{0,1 s}
{0,9 s}
\loigiai{
\begin{itemize}
\item Độ cứng k lò xo không thay đổi, từ $T=2\pi \sqrt{\dfrac{m}{k}}\Rightarrow T\sim \sqrt{m}$.
\item $\begin{cases}
T_1 \sim \sqrt{m_1}\\
T_2 \sim \sqrt {m_1} \\
T \sim \sqrt{m_1+m_2}\end{cases} 
\Rightarrow \begin{cases}
T_1^2 \sim m_1\\T_2^2 \sim m_2\\
T^2 \sim m_1 + m_2
\end{cases} \Rightarrow T^2 = T_1^2 + T_2^2
\Rightarrow T_1=\sqrt{T^2-T_2^2}=0,3$ s.
\end{itemize}
}
%\haidong{7}
\end{ex} \haidong{20}

\begin{ex}%[3P1K3-2][Loai1] 
Một lò xo nhẹ lần lượt liên kết với các vật có khối lượng $m_1$, $m_2$ và $m$ thì chu kì dao động lần lượt bằng $T_1=1,6\ s$, $T_2=1,8\ s$ và T. Nếu $m^2=2m_1^2+5m_2^2$ thì T bằng
\choice
{1,2 s}
{2,7 s}
{\True 2,8 s}
{4,6 s}
\loigiai{
T tỉ lệ thuận với $\sqrt{m}$ hay $m^2$ tỉ lệ với $T^4$ nên từ hệ thức $m^2=2m_1^2+5m_2^2$ suy ra: 
\begin{align*}
T^4=2T_1^4+5T_2^4\Rightarrow T=\sqrt[4]{2T_1^4+5T_2^4}\approx 2,8\ s.
\end{align*}}
%\haidong{5}
\end{ex} \haidong{20}
%\loai{Liên hệ các độ cứng và chu kì}
\begin{ex}%[3P1K3-2]
Một vật có khối lượng m treo vào một lò xo độ cứng $k_1$ thì chu kì dao động là $T_1 = 2$ s. Thay bằng lò xo có độ cứng $k_2$ thì chu kì dao động là $T_2 = 1,8$ s. Chu kì dao động của con lắc khi thay bằng một lò xo khác có độ cứng $k=3k_1+2k_2$ là
\choice
{0,73 s}
{\True 0,86 s}
{1,37 s}
{1,17 s}
\loigiai{
\begin{itemize}
\item Vật nặng khối lượng m không thay đổi, từ $T=2\pi \sqrt{\dfrac{m}{k}} \Rightarrow T\sim \dfrac{1}{\sqrt{k}}$.
\item $\begin{cases}
T_1 \sim \sqrt{\dfrac{1}{k_1}} \\
T_2 \sim \sqrt{\dfrac{1}{k_2}} \\
T \sim \sqrt{\dfrac{1}{3k_1 + 2k_2}} 
\end{cases} \Rightarrow 
\begin{cases}
\dfrac{1}{T_1^2} \sim k_1\\
\dfrac{1}{T_2^2} \sim k_2\\
\dfrac{1}{T^2} \sim 3k_1 + 2k_2
\end{cases} \Rightarrow 
\begin{cases}
\dfrac{3}{T_1^2} \sim 3k_1\\
\dfrac{2}{T_2^2} \sim 2k_2\\
\dfrac{1}{T^2} \sim 3k_1 + 2k_2
\end{cases} 
\Rightarrow \dfrac{1}{T^2} = \dfrac{3}{T_1^2} + \dfrac{2}{T_2^2} \Rightarrow T = 0,86\ s$.
\end{itemize}
}
%\haidong{5}
\end{ex} \haidong{20}

%\loai{Liên hệ chu kì và chiều dài lò xo}
\begin{ex}%[3P1G6-2][Loai1] 
Một con lắc lò xo được cấu tạo bởi một lò xo đồng nhất gắn với một vật nhỏ có khối lượng xác định. Chu kì dao động riêng của con lắc là 1,2 s. Nếu thay đổi độ dài lò xo đi 10 cm thì chu kì dao động riêng của con lắc là 0,8 s. Biết độ cứng k tỉ lệ nghịch với chiều dài tự nhiên của lò xo. Độ dài ban đầu của lò xo là
\choice
{15 cm}
{0 cm}
{40 cm}
{\True 18 cm}
\loigiai{
\begin{itemize}
\item Gọi $k_0$ là độ cứng của 1 m lò xo, $\ell$ là chiều dài ban đầu của lò xo.
\item Ta có: $\begin{cases}
T_1 = \dfrac{2\pi}{\omega} = 2\pi\sqrt{\dfrac{m\ell}{k_0}}\\
T_2 = 2\pi\sqrt{\dfrac{m\left(\ell-0,1\right)}{k_0}}
\end{cases}
\Rightarrow \dfrac{T_1}{T_2} = \sqrt{\dfrac{\ell}{\ell - 0,1}} = \dfrac{1,2}{0,8}
\Rightarrow \ell = 0,18\ m = 18\ cm$.
\end{itemize}
}
%\haidong{10}
\end{ex} \haidong{20}

\begin{ex}%[3P1G3-2][Loai1] 
Một lò xo đồng chất, tiết diện đều được cắt thành ba lò xo có chiều dài tự nhiên là $\ell$; $\ell-5$ và $\ell-8$ (đơn vị của $\ell$ là cm). Lần lượt gắn mỗi lò xo này với vật nhỏ khối lượng m thì được ba con lắc dao động với chu kì tương ứng là 2 s; $\sqrt{2}\ s$ và T. Biết độ cứng các lò xo tỉ lệ nghịch với chiều dài tự nhiên của nó. Giá trị của T là
\choice
{1,00 s}
{\True 0,89 s}
{1,41 s}
{0,50 s}
\loigiai{
\begin{itemize}
\item Gọi $k_0$ là độ cứng của 1 m lò xo.
\item Ta có: 
\begin{align*}
&T_1 = \dfrac{2\pi}{\omega} = 2\pi\sqrt{\dfrac{m\ell}{100k_0}} = 2 \quad(1)\\
&T_2 = 2\pi\sqrt{\dfrac{m\left(\ell-5\right)}{100k_0}} = \sqrt{2} \quad(2)\\
&T=2\pi\sqrt{\dfrac{m\left(\ell-8\right)}{100k_0}} \quad (3)
\end{align*}
\item Từ $(1)$ và $(2) \Rightarrow \dfrac{T_1}{T_2} = \sqrt{\dfrac{\ell}{\ell - 5}} = \sqrt{2} \Rightarrow \ell = 10$ cm.
\item Từ $(2)$ và $(3) \Rightarrow \dfrac{T_2}{T} = \sqrt{\dfrac{\ell-5}{\ell-8}} = \sqrt{\dfrac{5}{2}} \Rightarrow T = \dfrac{T_2}{\sqrt{\dfrac{5}{2}}} = \dfrac{\sqrt{2}}{\sqrt{\dfrac{5}{2}}} = 0,89\ s$.
\end{itemize}
}%<MyLT1>
%\haidong{10}
\end{ex} \haidong{20}

  \setcounter{ex}{0}
\PhanII
\begin{ex}
Xét một con lắc lò xo. Vật nặng có khối lượng m và lò xo có độ cứng k.
\choiceTF[t]
{\True Khi tăng khối lượng vật nặng 4 lần, chu kì tăng 2 lần}
{Chu kì tỉ lệ thuận với độ cứng của lò xo}
{\True Tần số tỉ lệ nghịch với căn bậc hai của khối lượng}
{Giảm $k$ đi 4 lần làm chu kì giảm một nửa}
\loigiai{
\begin{itemchoice}
\itemch 
\itemch 
\itemch 
\itemch 
\end{itemchoice}
}
\end{ex}

\begin{ex}
Một con lắc lò xo dao động với chu kì $T_0$. Thay lò xo khác có độ cứng gấp đôi và vẫn giữ nguyên khối lượng.
\choiceTF[t]
{Chu kì mới bằng $2T_0$}
{\True Tần số mới bằng $\sqrt{2}$ lần tần số cũ}
{\True Chu kì mới bằng $\dfrac{T_0}{\sqrt{2}}$}
{Tần số mới gấp đôi tần số cũ}
\loigiai{
\begin{itemchoice}
\itemch 
\itemch 
\itemch 
\itemch 
\end{itemchoice}
}
\end{ex}

\begin{ex}
Một vật khối lượng $m$ được treo vào lò xo và dao động với chu kì $T$, tần số $f$. Gắn thêm một vật khối lượng $m$ giống hệt vào vật cũ.
\choiceTF[t]
{\True Chu kì mới bằng $T\sqrt{2}$ \vphantom{$\dfrac{f}{\sqrt{2}}$}}
{Tần số mới bằng $2f$}
{\True Tần số mới bằng $\dfrac{f}{\sqrt{2}}$}
{Chu kì tăng $25\%$}
\loigiai{
\begin{itemchoice}
\itemch 
\itemch 
\itemch 
\itemch 
\end{itemchoice}
}
\end{ex}

\begin{ex}
Một lò xo đồng chất, tiết diện đều có độ cứng $k$. Cắt lò xo thành hai phần bằng nhau và gắn mỗi phần đó lần lượt với cùng một khối lượng $m$.
\choiceTF[t]
{\True Độ cứng của mỗi nửa lò xo bằng $2k$}
{Chu kì dao động của hệ mới gấp đôi chu kì ban đầu}
{\True Tần số dao động của hệ mới gấp $\sqrt{2}$ lần ban đầu}
{\True Chu kì dao động của hệ mới bằng $\dfrac{T}{\sqrt{2}}$}
\loigiai{
\begin{itemchoice}
\itemch 
\itemch 
\itemch 
\itemch 
\end{itemchoice}
}
\end{ex}

\setcounter{ex}{0}
\PhanIII
\begin{ex}
Một con lắc lò xo có độ cứng $k = 80\ N/m$, khối lượng $m = 0,20\ kg$. Chu kì dao động của con lắc là bao nhiêu giây (làm tròn kết quả đến chữ số hàng phần trăm)?
\shortans[oly]{0,32}
\loigiai{}
\end{ex}

\begin{ex}
Một con lắc lò xo dao động với chu kì $T = 0,4\ s$, lò xo có độ cứng $k = 50\ N/m$. Khối lượng của vật nặng bằng bao nhiêu gam (làm tròn kết quả đến chữ số hàng đơn vị)?
\shortans[oly]{127}

\end{ex}

\begin{ex}
Con lắc lò xo đang dao động với khối lượng $m$ và chu kì $1,5\ s$. Nếu giảm khối lượng đi $0,75\ kg$ thì chu kì còn $1\ s$. Khối lượng ban đầu $m$ bằng bao nhiêu kg?
\shortans[oly]{1,5}
\loigiai{}
\end{ex}

\begin{ex}
Một con lắc lò xo thực hiện $90$ dao động trong $30\ s$. Tần số dao động của con lắc là bao nhiêu Hz?
\shortans[oly]{3}
\loigiai{}
\end{ex}

\begin{ex}
Lò xo đồng chất dài tự nhiên $\ell_0 = 0,20\ m$, độ cứng $k_0$. Nếu kéo dãn lò xo thành $0,25\ m$ rồi gắn vào cùng một khối lượng $m$, chu kì mới đo được là $1,1\ s$. Chu kì ban đầu (với $\ell_0$) bằng bao nhiêu giây (làm tròn kết quả đến chữ số hàng phần trăm)?
\shortans[oly]{0,98}
\loigiai{}
\end{ex}

\begin{ex}
Một con lắc lò xo có chu kì $0,6\ s$. Hỏi trong $5$ phút con lắc thực hiện bao nhiêu dao động toàn phần?
\shortans[oly]{500}
\loigiai{}
\end{ex}

